% ===========================================================================
\section{Full UAIR$^+$ arithmetization}\label{sec:uair}
% ===========================================================================

We assemble the building blocks of \Cref{sec:arith} into a single
UAIR$^+$ instance. As in the SHA UAIR, we use the \emph{shift-access}
extension of the main paper (Remark~8.2 of the Zinc+ paper); we do
\emph{not} need the affine-combination lookup extension because the
ECDSA UAIR has no lookup constraints.

\subsection{Trace shape}\label{sec:trace-shape}

The trace runs $\mathrm{NUM\_SHAMIR\_ROUNDS} := 256$ Shamir
double-and-add iterations followed by one final row that holds the
output Jacobian state. Concretely, with
$\mathrm{FINAL\_ROW} := \mathrm{NUM\_SHAMIR\_ROUNDS} = 256$, the
active row range is $t \in [0, 256]$ ($257$ active rows total) and the
trace length is the smallest power of $2$ with
$2^{\nu} \ge 257$, namely $n := 2^{9} = 512$.

Per row $t$, the trace processes one Shamir step
$P[t+1] = 2P[t] + T[t]$ (with $T[t] \in \{\OOO, G, Q, G+Q\}$ chosen by
the bit pair $(b_1[t], b_2[t])$); the final row $256$ holds the
result $R$ for the off-protocol affine readout.

\subsection{Column layout}\label{sec:cols}

Following Definition~8.1 of the main paper, our index has
$\mathbf{q} = (p)$ where $p$ is the $\secp$ base prime, so we have a
single $\F_p$ branch. \emph{All} witness and constraint columns live
in this branch; there is no $\Q[X]$ branch and no binary-polynomial
section. The branch index $d_p = 1$ (treating $\F_p[X]$ at degree $0$)
makes every cell semantically an $\F_p$ element.

\emph{Storage vs.\ semantics.} The implementation does not — and
cannot — store $\F_p$ cells natively: in the framework, all branches
of the $\mathbf{q}$-tuple are realised as \emph{integer} columns
(\texttt{Int<5>}, i.e.\ $320$-bit fixed-precision integers; see
\texttt{EC\_FP\_INT\_LIMBS}~$=5$). Each cell holds the canonical
representative of its $\F_p$ value in $[0, p) \subset \Z$, and every
constraint of \Cref{sec:constraints} is an integer expression that
must lie in the principal ideal $(p) \subset \Z$. Constraints are written modulo $p$. In all constraints, the witness and public input are implicitly projected via the modular reduction map $\phi_p: \mathbb{Z}\to \mathbb{F}_p$.



Tables \ref{tab:pub-cols} and~\ref{tab:wit-cols} list the $15$ public
and $8$ witness columns, respectively. We write
$P[t]$ for the entry at row $t$ of column $P$ and
$P^{\downarrow}[t] := P[t+1]$ for shifted access, with $P^{\downarrow}[256]=0$.

\begin{table}[h]
\centering\small
\begin{tabular}{ll}
\toprule
\textbf{Column} & \textbf{Role} \\
\midrule
\multicolumn{2}{l}{\textit{Public columns ($15$ — entire ECDSA UAIR is public-input--gated)}} \\[1pt]
$\col{S\_INIT}$        & Selector: $1$ at row $0$, else $0$ \\
$\col{S\_ACTIVE}$      & Selector: $1$ for $t \in [0, 256)$, else $0$ \\
$\col{S\_FINAL}$       & Selector: $1$ at $\mathrm{FINAL\_ROW} = 256$, else $0$ \\
$\col{S\_ADD}$         & $\col{S\_ADD}[t] = b_1[t] + b_2[t] - b_1[t] b_2[t]$ \\
$\col{PA\_B1}$         & First scalar bit $b_1[t]$ of the row's Shamir bit pair (\Cref{rem:addend-cost}) \\
$\col{PA\_B2}$         & Second scalar bit $b_2[t]$ of the row's Shamir bit pair \\
$\col{PA\_QX}$         & Affine $x$-coord. of the public-key point $Q$ (constant across rows) \\
$\col{PA\_QY}$         & Affine $y$-coord.\ of $Q$ \\
$\col{PA\_QGX}$        & Affine $x$-coord. of $G + Q$ \\
$\col{PA\_QGY}$        & Affine $y$-coord. of $G + Q$ \\
$\col{PA\_R\_INIT\_X}$ & Initial Jacobian $X$-coordinate of $R_{\mathrm{init}}$ at row $0$ \\
$\col{PA\_R\_INIT\_Y}$ & Initial Jacobian $Y$-coordinate of $R_{\mathrm{init}}$ at row $0$ \\
$\col{PA\_R\_INIT\_Z}$ & Initial Jacobian $Z$-coordinate of $R_{\mathrm{init}}$ at row $0$ \\
$\col{PA\_Z\_INV}$     & $P_Z[256]^{-1} \bmod p$ at row $256$ (\Cref{con:f1}; cells at $t \neq 256$ are $0$) \\
$\col{PA\_R\_X}$       & Affine $x$-coordinate of $R$ at row $256$ (\Cref{con:f2}) \\
\bottomrule
\end{tabular}
\caption{Public columns ($15$ total). All are integer typed (i.e.\ \texttt{Int<5>}-stored, see \Cref{sec:cols}); the four selectors plus the bit pair are boolean. The boolean / constant-across-rows structural properties are intended to be discharged by verifier-side direct inspection of \texttt{public\_trace} (see soundness obligation in \Cref{rem:addend-cost}).}
\label{tab:pub-cols}
\end{table}

\begin{table}[h]
\centering\small
\begin{tabular}{ll}
\toprule
\textbf{Column} & \textbf{Role} \\
\midrule
\multicolumn{2}{l}{\textit{Witness columns ($8$ — all integer elements, no lookup constraints)}} \\[1pt]
$P_X$  & Chained Jacobian $X$-coordinate of $P[t]$ \quad (input at row $t$, output via $P_X^{\downarrow}$) \\
$P_Y$  & Chained Jacobian $Y$-coordinate of $P[t]$ \\
$P_Z$  & Chained Jacobian $Z$-coordinate of $P[t]$ \\
$P_{\mathrm{pa},X}$ & Doubled-point $X$-coordinate $X(2 P[t])$ \\
$P_{\mathrm{pa},Y}$ & Doubled-point $Y$-coordinate $Y(2 P[t])$ \\
$P_{\mathrm{pa},Z}$ & Doubled-point $Z$-coordinate $Z(2 P[t])$ \\
$H$    & Mixed-addition scratch $H = T_X \cdot P_{\mathrm{pa},Z}^2 - P_{\mathrm{pa},X}$ \eqref{eq:H} \quad (named $\col{W\_C}$ in the source) \\
$R_a$  & Mixed-addition scratch $R_a = T_Y \cdot P_{\mathrm{pa},Z}^3 - P_{\mathrm{pa},Y}$ \eqref{eq:Ra} \quad (named $\col{W\_D}$ in the source) \\
\bottomrule
\end{tabular}
\caption{Witness columns ($8$ total).}
\label{tab:wit-cols}
\end{table}

\begin{remark}[No range checks]\label{rem:no-range}
Although every cell is stored as an \texttt{Int<5>} integer
(\Cref{sec:cols}), no range-check or lookup constraint over the
<<<<<<< Updated upstream
trace columns is required. The reason is that onstraints encode $\F_p$ equalities, so non-canonical
rational representatives of the witness elements satisfy
=======
trace columns is required. The reason is twofold:
\emph{(i)}~constraints encode $\F_p$ equalities, so non-canonical
rational representatives of the wi ($x = x' + kp$ for any $k$) satisfy
>>>>>>> Stashed changes
the constraints exactly when their canonical reductions do — the
prover gains no leverage from picking non-canonical representatives. Further, if the prover uses rationals not in the localization ring $\mathbb{Z}_{(p)}$ (i.e.\ not reducible modulo $p$) then the verifier detects it with high probability, as per Theorem 5.4 of the Zinc+ paper.
\end{remark}

\begin{remark}[Inlined intermediates]\label{rem:inlined}
The higher-degree intermediates $S = P_Y^2$, $P_Y^4$, $P_{\mathrm{pa},Z}^2$,
$P_{\mathrm{pa},Z}^3$, $H^2$, $H^3$, and $P_{\mathrm{pa},X} H^2$ are
\emph{not} stored as committed columns; they appear inlined inside
the constraint expressions. Likewise the addition outputs
$P_{\mathrm{add}, X/Y/Z}$ from \eqref{eq:add-z}--\eqref{eq:add-y} are
not committed — they appear inlined inside the three
output-selection constraints \eqref{eq:next-x}--\eqref{eq:next-z}.
This trades higher per-constraint algebraic degree for fewer
committed columns; the maximum constraint degree in the trace
MLEs is~$6$.
\end{remark}

\begin{remark}[Conditional add via output-selection rather than gating]
\label{rem:no-add-witness}
The per-row addition output is not committed; the conditional add
$P_{\mathrm{add}}$-or-$P_{\mathrm{pa}}$ is folded into the
chaining via \eqref{eq:next-x}--\eqref{eq:next-z}. This avoids
the dedicated $T = (T_X : T_Y : T_Z)$ Jacobian-coordinate columns that an
explicit \emph{compute-then-gate} approach would require.
\end{remark}

\subsection{Constraints}\label{sec:constraints}

There are $11$ constraints in total: $3$ doubling constraints, $2$
mixed-addition scratch constraints, $3$ output-selection constraints
(which fold in the conditional add and the next-row chaining), and
$3$ init-boundary constraints. All constraints carry the zero ideal,
i.e.\ they are strict equalities over $\F_p$.

\paragraph{Doubling constraints (anchored at every active row $t$).}

For every $t \in [0, 256)$ (gated by $\col{S\_ACTIVE}$):
\begin{align}
\col{S\_ACTIVE}[t] \cdot \bigl(P_{\mathrm{pa},Z}[t] \;-\; 2\, P_Y[t]\, P_Z[t]\bigr) &= 0,\tag{D2}\label{con:d2}\\
\col{S\_ACTIVE}[t] \cdot \bigl(P_{\mathrm{pa},X}[t] \;-\; 9\, P_X[t]^4 \;+\; 8\, P_X[t]\, P_Y[t]^2\bigr) &= 0,\tag{D3}\label{con:d3}\\
\col{S\_ACTIVE}[t] \cdot \bigl(P_{\mathrm{pa},Y}[t] \;-\; 12\, P_X[t]^3\, P_Y[t]^2 \;+\; 3\, P_X[t]^2\, P_{\mathrm{pa},X}[t] \;+\; 8\, P_Y[t]^4\bigr) &= 0.\tag{D4}\label{con:d4}
\end{align}
Algebraic degrees in the trace MLEs (after gating): \ref{con:d2}~deg~$3$, \ref{con:d3}~deg~$5$, \ref{con:d4}~deg~$6$ (via the $12 P_X^3 P_Y^2$ term).

\paragraph{Addition-scratch constraints.}

For every $t \in [0, 256)$, with the in-circuit affine addend
$T_x, T_y$ from \eqref{eq:addend-formula-expanded} substituted
inline:
\begin{align}
\col{S\_ACTIVE}[t] \cdot \bigl(H[t] \;-\; T_x[t] \cdot P_{\mathrm{pa},Z}[t]^2 \;+\; P_{\mathrm{pa},X}[t]\bigr) &= 0,\tag{A1}\label{con:a1}\\
\col{S\_ACTIVE}[t] \cdot \bigl(R_a[t] \;-\; T_y[t] \cdot P_{\mathrm{pa},Z}[t]^3 \;+\; P_{\mathrm{pa},Y}[t]\bigr) &= 0,\tag{A2}\label{con:a2}
\end{align}
where
\[
  T_\square[t] \;=\; b_1[t]\, G_\square \;+\; b_2[t]\, Q_\square[t]
                  \;+\; b_1[t]\, b_2[t]\,\bigl((G+Q)_\square[t] - G_\square - Q_\square[t]\bigr)
\]
for $\square \in \{x, y\}$, with $b_1[t] = \col{PA\_B1}[t]$,
$b_2[t] = \col{PA\_B2}[t]$, and the column / scalar substitutions
$Q_x = \col{PA\_QX}, Q_y = \col{PA\_QY}, (G+Q)_x = \col{PA\_QGX},
(G+Q)_y = \col{PA\_QGY}$, with $G_x, G_y$ as UAIR scalars.
Algebraic degrees (after gating): $T_\square$ alone is degree $3$
in trace cells (the $b_1 b_2 \cdot Q_\square$ term reaches three
columns); together with $P_{\mathrm{pa},Z}^k$ and $\col{S\_ACTIVE}$,
\ref{con:a1}~deg~$6$, \ref{con:a2}~deg~$7$.

\paragraph{Output-selection-and-chaining constraints.}

These three constraints inline the addition outputs of
\eqref{eq:add-z}--\eqref{eq:add-y} into the conditional-add chaining
\eqref{eq:next-x}--\eqref{eq:next-z}. For every $t \in [0, 256)$:
\begin{align}
\col{S\_ACTIVE}[t] \cdot \Bigl(P_X^{\downarrow}[t] - P_{\mathrm{pa},X}[t]
   - \col{S\_ADD}[t] \bigl(R_a[t]^2 - H[t]^3 - 2 P_{\mathrm{pa},X}[t] H[t]^2 - P_{\mathrm{pa},X}[t]\bigr)\Bigr) &= 0, \tag{O1}\label{con:o1}\\[4pt]
\col{S\_ACTIVE}[t] \cdot \Bigl(P_Y^{\downarrow}[t] - P_{\mathrm{pa},Y}[t]
   - \col{S\_ADD}[t] \bigl(3 R_a[t] P_{\mathrm{pa},X}[t] H[t]^2 + R_a[t] H[t]^3 \notag\\
   {} - R_a[t]^3 - P_{\mathrm{pa},Y}[t] H[t]^3 - P_{\mathrm{pa},Y}[t]\bigr)\Bigr) &= 0, \tag{O2}\label{con:o2}\\[4pt]
\col{S\_ACTIVE}[t] \cdot \Bigl(P_Z^{\downarrow}[t] - P_{\mathrm{pa},Z}[t]
   - \col{S\_ADD}[t] \bigl(P_{\mathrm{pa},Z}[t] H[t] - P_{\mathrm{pa},Z}[t]\bigr)\Bigr) &= 0. \tag{O3}\label{con:o3}
\end{align}
Algebraic degrees (after gating): \ref{con:o1}~deg~$5$, \ref{con:o2}~deg~$6$, \ref{con:o3}~deg~$4$. The deg-$6$ monomial in
\ref{con:o2} is exactly $\col{S\_ACTIVE} \cdot \col{S\_ADD} \cdot R_a \cdot P_{\mathrm{pa},X} \cdot H^2$ (\Cref{rem:inlined-add}).

\paragraph{Init-boundary constraints (anchored at row $0$).}

The Jacobian state at row $0$ is pinned to the verifier-supplied
seed $R_{\mathrm{init}}$:
\begin{align}
\col{S\_INIT}[t] \cdot \bigl(P_X[t] - \col{PA\_R\_INIT\_X}[t]\bigr) &= 0,\tag{I1}\label{con:init-x}\\
\col{S\_INIT}[t] \cdot \bigl(P_Y[t] - \col{PA\_R\_INIT\_Y}[t]\bigr) &= 0,\tag{I2}\label{con:init-y}\\
\col{S\_INIT}[t] \cdot \bigl(P_Z[t] - \col{PA\_R\_INIT\_Z}[t]\bigr) &= 0.\tag{I3}\label{con:init-z}
\end{align}
Algebraic degrees (after gating): all three are deg~$2$.

\paragraph{Final-row affine-readout constraints (anchored at row $\mathrm{FINAL\_ROW} = 256$).}

The two $\col{S\_FINAL}$-gated $(p)$-ideal constraints
\ref{con:f1}, \ref{con:f2} of \Cref{sec:boundary} pin the
public columns $\col{PA\_Z\_INV}$ and $\col{PA\_R\_X}$ to the affine
$x$-readout of the final-row Jacobian state:
\begin{align}
\col{S\_FINAL}[t] \cdot \bigl(P_Z[t] \cdot \col{PA\_Z\_INV}[t] - 1\bigr) &\in (p), \tag{F1}\\
\col{S\_FINAL}[t] \cdot \bigl(P_X[t] \cdot \col{PA\_Z\_INV}[t]^2 - \col{PA\_R\_X}[t]\bigr) &\in (p). \tag{F2}
\end{align}
Algebraic degrees (after gating): \ref{con:f1}~deg~$3$
($\col{S\_FINAL} \cdot P_Z \cdot \col{PA\_Z\_INV}$),
\ref{con:f2}~deg~$4$ ($\col{S\_FINAL} \cdot P_X \cdot \col{PA\_Z\_INV}^2$).
The order-field equality $\col{PA\_R\_X}[256] \equiv r \pmod n$ is
checked off-protocol by the verifier (\Cref{rem:final-readout}); no
in-circuit constraint is needed for it.

\subsection{Algebraic-degree summary}\label{sec:degree}

\begin{center}
\small
\begin{tabular}{lll}
\toprule
\textbf{Constraint} & \textbf{Max.\ deg.\ (after gating)} & \textbf{Reason} \\
\midrule
\ref{con:d2}     & $3$ & $\col{S\_ACTIVE} \cdot P_Y \cdot P_Z$ \\
\ref{con:d3}     & $5$ & $\col{S\_ACTIVE} \cdot P_X^4$ \\
\ref{con:d4}     & $6$ & $\col{S\_ACTIVE} \cdot P_X^3 \cdot P_Y^2$ \\
\ref{con:a1}     & $6$ & $\col{S\_ACTIVE} \cdot \col{PA\_B1} \cdot \col{PA\_B2} \cdot \col{PA\_QGX} \cdot P_{\mathrm{pa},Z}^2$ \\
\ref{con:a2}     & $7$ & $\col{S\_ACTIVE} \cdot \col{PA\_B1} \cdot \col{PA\_B2} \cdot \col{PA\_QGY} \cdot P_{\mathrm{pa},Z}^3$ \\
\ref{con:o1}     & $5$ & $\col{S\_ACTIVE} \cdot \col{S\_ADD} \cdot P_{\mathrm{pa},X} \cdot H^2$ \\
\ref{con:o2}     & $6$ & $\col{S\_ACTIVE} \cdot \col{S\_ADD} \cdot R_a \cdot P_{\mathrm{pa},X} \cdot H^2$ \\
\ref{con:o3}     & $4$ & $\col{S\_ACTIVE} \cdot \col{S\_ADD} \cdot P_{\mathrm{pa},Z} \cdot H$ \\
\ref{con:init-x}--\ref{con:init-z} & $2$ & $\col{S\_INIT} \cdot P_{X/Y/Z}$ \\
\ref{con:f1}     & $3$ & $\col{S\_FINAL} \cdot P_Z \cdot \col{PA\_Z\_INV}$ \\
\ref{con:f2}     & $4$ & $\col{S\_FINAL} \cdot P_X \cdot \col{PA\_Z\_INV}^2$ \\
\midrule
\textbf{Global maximum} & \textbf{$7$} & attained by \ref{con:a2}\footnote{In-circuit $T_y$ is degree $3$, multiplied by $P_{\mathrm{pa},Z}^3$ deg $3$ and $\col{S\_ACTIVE}$} \\
\bottomrule
\end{tabular}
\end{center}
